\begin{ex}
Tính: $14\text{ cm} + 5\text{ cm}$. Kết quả là bao nhiêu xăng-ti-mét?
\shortans{19}
\loigiai{
Ta có: $14\text{ cm} + 5\text{ cm} = 19\text{ cm}$.
\textbf{Đáp số:} $19$.
}
\end{ex}

\begin{ex}
Tính: $18\text{ cm} - 6\text{ cm}$. Kết quả là bao nhiêu xăng-ti-mét?
\shortans{12}
\loigiai{
Ta có: $18\text{ cm} - 6\text{ cm} = 12\text{ cm}$.
\textbf{Đáp số:} $12$.
}
\end{ex}

\begin{ex}
Tính: $20\text{ cm} + 30\text{ cm}$. Kết quả là bao nhiêu xăng-ti-mét?
\shortans{50}
\loigiai{
Ta có: $20\text{ cm} + 30\text{ cm} = 50\text{ cm}$.
\textbf{Đáp số:} $50$.
}
\end{ex}

\begin{ex}
Tính: $70\text{ cm} - 40\text{ cm}$. Kết quả là bao nhiêu xăng-ti-mét?
\shortans{30}
\loigiai{
Ta có: $70\text{ cm} - 40\text{ cm} = 30\text{ cm}$.
\textbf{Đáp số:} $30$.
}
\end{ex}

\begin{ex}
Điền số thích hợp: $12\text{ cm} + \dots\text{ cm} = 18\text{ cm}$.
\shortans{6}
\loigiai{
Số cần điền là: $18 - 12 = 6$.
\textbf{Đáp số:} $6$.
}
\end{ex}

\begin{ex}
Điền số thích hợp: $19\text{ cm} - \dots\text{ cm} = 15\text{ cm}$.
\shortans{4}
\loigiai{
Số cần điền là: $19 - 15 = 4$.
\textbf{Đáp số:} $4$.
}
\end{ex}

\begin{ex}
Đoạn thẳng AB dài 11 cm, đoạn thẳng CD dài 8 cm. Hỏi cả hai đoạn thẳng dài tất cả bao nhiêu xăng-ti-mét?
\shortans{19}
\loigiai{
Tổng độ dài hai đoạn thẳng là: $11 + 8 = 19$ (cm).
\textbf{Đáp số:} $19$.
}
\end{ex}

\begin{ex}
Đoạn thẳng AB dài hơn đoạn thẳng CD ở bài toán trên bao nhiêu xăng-ti-mét?
\shortans{3}
\loigiai{
AB dài hơn CD là: $11 - 8 = 3$ (cm).
\textbf{Đáp số:} $3$.
}
\end{ex}

\begin{ex}
Một sợi dây thừng dài 50 cm. Bác thợ cắt bớt đi 20 cm. Hỏi sợi dây thừng còn lại dài bao nhiêu xăng-ti-mét?
\shortans{30}
\loigiai{
Chiều dài sợi dây thừng còn lại là: $50 - 20 = 30$ (cm).
\textbf{Đáp số:} $30$.
}
\end{ex}

\begin{ex}
Tính giá trị biểu thức: $15\text{ cm} + 4\text{ cm} - 9\text{ cm}$.
\shortans{10}
\loigiai{
Thực hiện lần lượt: $15 + 4 = 19$; $19 - 9 = 10$ (cm).
\textbf{Đáp số:} $10$.
}
\end{ex}

\begin{ex}
Tính giá trị biểu thức: $18\text{ cm} - 8\text{ cm} + 5\text{ cm}$.
\shortans{15}
\loigiai{
Thực hiện lần lượt: $18 - 8 = 10$; $10 + 5 = 15$ (cm).
\textbf{Đáp số:} $15$.
}
\end{ex}

\begin{ex}
Một cuộn ruy băng dài 40 cm, mẹ mua thêm một cuộn ruy băng dài 30 cm. Hỏi mẹ có tất cả bao nhiêu xăng-ti-mét ruy băng?
\shortans{70}
\loigiai{
Tổng số ruy băng mẹ có là: $40 + 30 = 70$ (cm).
\textbf{Đáp số:} $70$.
}
\end{ex}

\begin{ex}
Chiếc khăn quàng đỏ dài 80 cm, chiếc nơ trang trí dài 20 cm. Chiếc khăn quàng đỏ dài hơn chiếc nơ bao nhiêu xăng-ti-mét?
\shortans{60}
\loigiai{
Khăn quàng đỏ dài hơn nơ là: $80 - 20 = 60$ (cm).
\textbf{Đáp số:} $60$.
}
\end{ex}

\begin{ex}
Một thanh nhôm dài 16 cm. Người ta cắt thành hai đoạn bằng nhau. Hỏi mỗi đoạn dài bao nhiêu xăng-ti-mét?
\shortans{8}
\loigiai{
Vì hai đoạn bằng nhau và $8 + 8 = 16$ nên mỗi đoạn dài $8$ cm.
\textbf{Đáp số:} $8$.
}
\end{ex}

\begin{ex}
Đoạn thẳng MN dài 7 cm. Nối tiếp thêm đoạn thẳng NP dài 8 cm và đoạn PQ dài 5 cm (M, N, P, Q thẳng hàng). Hỏi đoạn MQ dài bao nhiêu xăng-ti-mét?
\shortans{20}
\loigiai{
Độ dài đoạn MQ là: $7 + 8 + 5 = 20$ (cm).
\textbf{Đáp số:} $20$.
}
\end{ex}

\begin{ex}
Một cây nến ban đầu cao 20 cm. Sau khi thắp sáng một thời gian, nến còn lại cao 12 cm. Hỏi cây nến đã bị cháy ngắn đi bao nhiêu xăng-ti-mét?
\shortans{8}
\loigiai{
Độ dài nến đã cháy là: $20 - 12 = 8$ (cm).
\textbf{Đáp số:} $8$.
}
\end{ex}

\begin{ex}
Một con kiến bò từ điểm A đến điểm B dài 9 cm, rồi bò tiếp từ B đến C dài 10 cm. Hỏi con kiến đã bò được tất cả bao nhiêu xăng-ti-mét?
\shortans{19}
\loigiai{
Tổng quãng đường kiến bò là: $9 + 10 = 19$ (cm).
\textbf{Đáp số:} $19$.
}
\end{ex}

\begin{ex}
Tìm số tự nhiên $x$ thỏa mãn: $x\text{ cm} + 5\text{ cm} = 17\text{ cm}$.
\shortans{12}
\loigiai{
Ta có: $x = 17 - 5 = 12$.
\textbf{Đáp số:} $12$.
}
\end{ex}

\begin{ex}
Tìm số tự nhiên $y$ thỏa mãn: $20\text{ cm} - y\text{ cm} = 14\text{ cm}$.
\shortans{6}
\loigiai{
Ta có: $y = 20 - 14 = 6$.
\textbf{Đáp số:} $6$.
}
\end{ex}

\begin{ex}
Một cái hộp hình chữ nhật có chiều dài 15 cm, chiều rộng 10 cm. Chiều dài hơn chiều rộng bao nhiêu xăng-ti-mét?
\shortans{5}
\loigiai{
Chiều dài hơn chiều rộng là: $15 - 10 = 5$ (cm).
\textbf{Đáp số:} $5$.
}
\end{ex}

\begin{ex}
Băng giấy màu xanh dài 25 cm, băng giấy màu đỏ dài 20 cm. Cả hai băng giấy dài bao nhiêu xăng-ti-mét?
\shortans{45}
\loigiai{
Tổng độ dài hai băng giấy là: $25 + 20 = 45$ (cm).
\textbf{Đáp số:} $45$.
}
\end{ex}

\begin{ex}
Một cây tre dài 90 cm. Chú thợ cắt bớt đi 30 cm để làm cán chổi. Cây tre còn lại dài bao nhiêu xăng-ti-mét?
\shortans{60}
\loigiai{
Cây tre còn lại dài: $90 - 30 = 60$ (cm).
\textbf{Đáp số:} $60$.
}
\end{ex}

\begin{ex}
Ba bạn kiến cùng bò trên một cành cây: kiến đen bò 10 cm, kiến vàng bò 12 cm, kiến lửa bò 8 cm. Hỏi kiến vàng bò nhiều hơn kiến lửa bao nhiêu xăng-ti-mét?
\shortans{4}
\loigiai{
Kiến vàng bò nhiều hơn kiến lửa là: $12 - 8 = 4$ (cm).
\textbf{Đáp số:} $4$.
}
\end{ex}

\begin{ex}
Tổng quãng đường bò của kiến đen và kiến lửa ở bài toán trên là bao nhiêu xăng-ti-mét?
\shortans{18}
\loigiai{
Tổng quãng đường là: $10 + 8 = 18$ (cm).
\textbf{Đáp số:} $18$.
}
\end{ex}

\begin{ex}
Một đoạn dây điện dài 60 cm, người ta dùng hết 40 cm để nối bóng đèn. Đoạn dây điện còn lại dài bao nhiêu xăng-ti-mét?
\shortans{20}
\loigiai{
Đoạn dây điện còn lại dài: $60 - 40 = 20$ (cm).
\textbf{Đáp số:} $20$.
}
\end{ex}
